%% kapa: exam=1 round=25 year=2014 period=1 form=A subject=economics subject_ko=경제원론 seq=20 number=60
%% source: Q-Net 감정평가사 기출문제 (cst00309) artlSeq=5101207 file=extracted/1008/5101207/2014년 제25회 감정평가사 1차시험 기출문제/1교시(A형).hwp
%% conversion: kordoc 4.24.51 -> tools/gen1.py (auto); review=proofread-sonnet; flags=none
\begin{question}{60}
효율적인 자원배분에 관한 설명으로 옳지 않은 것은? (단, $X$, $Y$ 두 재화 / A, B 2인 / $L$, $K$ 2 생산요소 경제를 가정한다.)
\begin{choices}
  \item 소비의 효율성 조건은 두 재화에 대한 $A$와 $B$의 한계대체율이 같을 때 만족한다.
  \item 생산의 효율성 조건은 두 생산요소에 대한 $X$와 $Y$의 한계기술대체율이 같을 때 만족한다.
  \item 생산의 효율성 조건을 만족하는 점들을 이어서 연결한 선을 생산의 계약곡선이라 한다.
  \item 생산의 효율성 조건을 만족하는 점들을 $X-Y$평면으로 옮겨 놓은 것을 생산가능곡선(production possibility curve)이라 한다.
  \item 생산의 효율성 조건, 소비의 효율성 조건, 종합적 효율성 조건을 모두 만족하는 점들을 $X-Y$평면으로 옮겨 놓은 것을 효용가능곡선(utility possibility curve)이라 한다.
\end{choices}
\end{question}
\answer{5}
